%======================================================================
% CAPÍTULO 18 - IMPLEMENTACIÓN
%======================================================================
\chapter{Implementación}

\bigskip

Este capítulo aborda los detalles técnicos de implementación, incluyendo el stack de tecnologías, consideraciones de despliegue y operación del sistema.

\bigskip

\section{Stack Tecnológico Detallado}

Sentinel está implementado sobre un stack tecnológico moderno que aprovecha las capacidades del ecosistema Rust:

\bigskip

\begin{center}
\begin{tabular}{lll}
\toprule
\textbf{Componente} & \textbf{Tecnología} & \textbf{Versión} \\
\midrule
Lenguaje & Rust & 2021 Edition \\
Runtime Asíncrono & Tokio & 1.x \\
Protocolo & gRPC & Protobuf v3 \\
Serialización & serde\_json & Latest \\
Paralelismo & Rayon & Latest \\
Scripting & Rhai & Latest \\
Fechas & chrono & Latest \\
\bottomrule
\end{tabular}
\end{center}

\bigskip

\section{Arquitectura de Código}

El código fuente está organizado en módulos:

\bigskip

\begin{center}
\begin{tabular}{llp{8cm}}
\toprule
\textbf{Módulo} & \textbf{Responsabilidad} \\
\midrule
kernel & Orquestación y coordinación general \\
calc & Motor de cálculo y fórmulas \\
model & Estructuras de datos \\
system & Configuración y utilidades \\
util & Funciones auxiliares \\
\bottomrule
\end{tabular}
\end{center}

\bigskip

\section{Deserialización Zero-Copy}

Sentinel utiliza serde\_json::from\_slice que permite la deserialización directa desde bytes sin copia adicional de memoria, aprovechando las instrucciones SIMD disponibles en los procesadores modernos.

\bigskip

\section{Pipeline Asíncrono}

El procesamiento sigue un patrón de pipeline asíncrono:

\begin{enumerate}
    \item El hilo principal recibe paquetes de red.
    \item Delega deserialización a un pool de hilos (tokio::spawn).
    \item Mientras la CPU procesa el lote N, la red descarga el lote N+1.
\end{enumerate}

\bigskip

\section{Archivos de Exportación}

\subsection{nomina\_exportada.csv}

Archivo principal con todos los campos de nómina:

\bigskip

Ejemplo de formato:

\begin{lstlisting}[breaklines=true,frame=none,basicstyle=\ttfamily\small]
cedula,nombres,apellidos,sexo,edo_civil,n_hijos,componente_id,grado_id,...
\end{lstlisting}

\bigskip

\subsection{aporte\_CICLO.csv}

Archivo de aportes para distribución de garantías:

\bigskip

Ejemplo de formato:

\begin{lstlisting}[breaklines=true,frame=none,basicstyle=\ttfamily\small]
cedula,nombres,apellidos,numero_cuenta,garantia_original,factor_aplicado,garantia_anticipo
10002142,,MATHEUS HIDALGO LUIS EDUARDO,01020451850000065650,615.90,0.70114191,431.83
\end{lstlisting}

\bigskip

\section{Consideraciones de Despliegue}

Requisitos mínimos:

\begin{itemize}
    \item 16 GB de RAM (recomendado 32 GB para 500k+ registros)
    \item Procesador multi-núcleo
    \item Conexión gRPC al servidor Sandra
    \item Espacio en disco para archivos CSV de salida
\end{itemize}

\bigskip

\section{Métricas de Rendimiento}

\begin{center}
\begin{tabular}{lll}
\toprule
\textbf{Métrica} & \textbf{Valor} & \textbf{Condiciones} \\
\midrule
Tiempo total & ~4 segundos & 500,000 registros \\
Throughput & ~125,000 reg/s & Procesamiento paralelo \\
Memoria & ~8 GB & 500k registros en RAM \\
\bottomrule
\end{tabular}
\end{center}

\vfill
\newpage
